% GENERATED FILE. Edit canonical lessons, then run npm run generate:books.
% canonical-source-hash: fnv1a64:055814b19e0c6419
% canonical-lessons: ZH-C135-zuqiu, ZH-C135-lanqiu, ZH-C135-ti, ZH-C135-bisai, ZH-C135-duanlian

\chapter{Sport}
\label{ch:zh-sport}
\hlchaptermodality{135}

\begin{chapteropening}
\textbf{By the end of this chapter:} \emph{I can talk about football, basketball, matches and working out.}

Complete the last lesson of chapter 135: I can talk about football, basketball, matches and working out.
\end{chapteropening}

\section[zúqiú]{zúqiú — football; soccer}

\label{lesson:ZH-C135-zuqiu}



\begin{quote}
Before the new one: say the Mandarin for empty, then the Mandarin for full.
\end{quote}

\subsection*{You'll want to know: zúqiú}
\emph{zúqiú} — \textquotedblleft{}football; soccer\textquotedblright{}.

\begin{quote}
\emph{wǒ xǐhuan zúqiú} — I love football.
\end{quote}

Its characters come later; for now it is a word for the ear.

\begin{cousinweb}
\emph{zúqiú} is built from \emph{zú} 'foot' and \emph{qiú} 'ball', both inherited from Middle Chinese, \emph{tsjowk} and \emph{gjuw} (Guangyun, in Baxter's transcription). The dictionary glosses \emph{gjuw} as 'fine jade', and under another way of writing it as a ball of fur that is struck.
\end{cousinweb}

\begin{grammarlens}[title={What you've built}]
One more word for what I like.
\end{grammarlens}

\subsection*{Your turn}
\begin{itemize}
\raggedright
  \item \emph{Say it:} \emph{zúqiú}
  \item \emph{Say it:} \emph{zúqiú}, once more
  \item \emph{Say it:} \emph{wǒ xǐhuan zúqiú} — I love football.
  \item \emph{Recall:} say the Mandarin for empty, then the Mandarin for full, then say \emph{zúqiú} again
\end{itemize}

\begin{tcolorbox}[breakable,colback=teal!4,colframe=teal!35!black,title={Before you move on}]
What does zúqiú mean? (\textquotedblleft{}Football; soccer\textquotedblright{}.) Say it once more.
\end{tcolorbox}
\section[lánqiú]{lánqiú — basketball}

\label{lesson:ZH-C135-lanqiu}



\begin{quote}
Before the new one: say the Mandarin for full, then the Mandarin for football.
\end{quote}

\subsection*{You'll want to know: lánqiú}
\emph{lánqiú} — \textquotedblleft{}basketball\textquotedblright{}.

\begin{quote}
\emph{wǒ gēge hěn gāo tā xǐhuan lánqiú} — My brother is tall, and he loves basketball.
\end{quote}

Its characters come later; for now it is a word for the ear.

\begin{cousinweb}
\emph{lánqiú} is built from \emph{lán} 'basket' and \emph{qiú} 'ball'. The first is inherited from Middle Chinese \emph{lam} (Guangyun, in Baxter's transcription), glossed 'a basket', a different word from the \emph{lán} of blue; the second is the \emph{qiú} of football.
\end{cousinweb}

\begin{grammarlens}[title={What you've built}]
One more word for what I like.
\end{grammarlens}

\subsection*{Your turn}
\begin{itemize}
\raggedright
  \item \emph{Say it:} \emph{lánqiú}
  \item \emph{Say it:} \emph{lánqiú}, once more
  \item \emph{Say it:} \emph{wǒ gēge hěn gāo tā xǐhuan lánqiú} — My brother is tall, and he loves basketball.
  \item \emph{Recall:} say the Mandarin for full, then the Mandarin for football, then say \emph{lánqiú} again
\end{itemize}

\begin{tcolorbox}[breakable,colback=teal!4,colframe=teal!35!black,title={Before you move on}]
What does lánqiú mean? (\textquotedblleft{}Basketball\textquotedblright{}.) Say it once more.
\end{tcolorbox}
\section[tī]{tī — to kick}

\label{lesson:ZH-C135-ti}



\begin{quote}
Before the new one: say the Mandarin for football, then the Mandarin for basketball.
\end{quote}

\subsection*{You'll want to know: tī}
\emph{tī} — \textquotedblleft{}to kick\textquotedblright{}.

\begin{quote}
\emph{wǒmen qù tī zúqiú} — We're off to play football.
\end{quote}

Its characters come later; for now it is a word for the ear.

\begin{cousinweb}
\emph{tī} The character is in the Guangyun as \emph{thek} (in Baxter's transcription), but only in the name of a beast; the verb is not given there. No dictionary we use records an older form, so its history is not settled.
\end{cousinweb}

\begin{grammarlens}[title={What you've built}]
One more word for what I like.
\end{grammarlens}

\subsection*{Your turn}
\begin{itemize}
\raggedright
  \item \emph{Say it:} \emph{tī}
  \item \emph{Say it:} \emph{tī}, once more
  \item \emph{Say it:} \emph{wǒmen qù tī zúqiú} — We're off to play football.
  \item \emph{Recall:} say the Mandarin for football, then the Mandarin for basketball, then say \emph{tī} again
\end{itemize}

\begin{tcolorbox}[breakable,colback=teal!4,colframe=teal!35!black,title={Before you move on}]
What does tī mean? (\textquotedblleft{}To kick\textquotedblright{}.) Say it once more.
\end{tcolorbox}
\section[bǐsài]{bǐsài — match; competition}

\label{lesson:ZH-C135-bisai}



\begin{quote}
Before the new one: say the Mandarin for basketball, then the Mandarin for to kick.
\end{quote}

\subsection*{You'll want to know: bǐsài}
\emph{bǐsài} — \textquotedblleft{}match; competition\textquotedblright{}.

\begin{quote}
\emph{jīntiān yǒu zúqiú bǐsài} — There's a football match today.
\end{quote}

Its characters come later; for now it is a word for the ear.

\begin{cousinweb}
\emph{bǐsài} is built from \emph{bǐ} 'to compare' and \emph{sài}, both inherited from Middle Chinese, \emph{pjijX} and \emph{sojH} (Guangyun, in Baxter's transcription). The dictionary glosses the second 'to repay'; how the pair came to mean a match is not settled.
\end{cousinweb}

\begin{grammarlens}[title={What you've built}]
One more word for what I like.
\end{grammarlens}

\subsection*{Your turn}
\begin{itemize}
\raggedright
  \item \emph{Say it:} \emph{bǐsài}
  \item \emph{Say it:} \emph{bǐsài}, once more
  \item \emph{Say it:} \emph{jīntiān yǒu zúqiú bǐsài} — There's a football match today.
  \item \emph{Recall:} say the Mandarin for basketball, then the Mandarin for to kick, then say \emph{bǐsài} again
\end{itemize}

\begin{tcolorbox}[breakable,colback=teal!4,colframe=teal!35!black,title={Before you move on}]
What does bǐsài mean? (\textquotedblleft{}Match; competition\textquotedblright{}.) Say it once more.
\end{tcolorbox}
\section[duànliàn]{duànliàn — to exercise; to work out}

\label{lesson:ZH-C135-duanlian}



\begin{quote}
Before the new one: say the Mandarin for to kick, then the Mandarin for a match.
\end{quote}

\subsection*{You'll want to know: duànliàn}
\emph{duànliàn} — \textquotedblleft{}to exercise; to work out\textquotedblright{}.

\begin{quote}
\emph{wǒ xǐhuan zǎoshang duànliàn} — I like working out in the morning.
\end{quote}

Its characters come later; for now it is a word for the ear.

\begin{cousinweb}
\emph{duànliàn} is built from \emph{duàn} 'to forge iron' and \emph{liàn} 'to refine metal', both inherited from Middle Chinese, \emph{twanH} and \emph{lenH} (Guangyun, in Baxter's transcription). The pair said of the body is not given there.
\end{cousinweb}

\begin{grammarlens}[title={What you've built}]
One more word for what I like.
\end{grammarlens}

\subsection*{Your turn}
\begin{itemize}
\raggedright
  \item \emph{Say it:} \emph{duànliàn}
  \item \emph{Say it:} \emph{duànliàn}, once more
  \item \emph{Say it:} \emph{wǒ xǐhuan zǎoshang duànliàn} — I like working out in the morning.
  \item \emph{Recall:} say the Mandarin for to kick, then the Mandarin for a match, then say \emph{duànliàn} again
\end{itemize}

\begin{tcolorbox}[breakable,colback=teal!4,colframe=teal!35!black,title={Before you move on}]
What does duànliàn mean? (\textquotedblleft{}To exercise; to work out\textquotedblright{}.) Say it once more.
\end{tcolorbox}
